\documentclass[]{article}
\usepackage{lmodern}
\usepackage{amssymb,amsmath}
\usepackage{ifxetex,ifluatex}
\usepackage{fixltx2e} % provides \textsubscript
\ifnum 0\ifxetex 1\fi\ifluatex 1\fi=0 % if pdftex
  \usepackage[T1]{fontenc}
  \usepackage[utf8]{inputenc}
\else % if luatex or xelatex
  \ifxetex
    \usepackage{mathspec}
  \else
    \usepackage{fontspec}
  \fi
  \defaultfontfeatures{Ligatures=TeX,Scale=MatchLowercase}
\fi
% use upquote if available, for straight quotes in verbatim environments
\IfFileExists{upquote.sty}{\usepackage{upquote}}{}
% use microtype if available
\IfFileExists{microtype.sty}{%
\usepackage{microtype}
\UseMicrotypeSet[protrusion]{basicmath} % disable protrusion for tt fonts
}{}
\usepackage[margin=1in]{geometry}
\usepackage{hyperref}
\hypersetup{unicode=true,
            pdfborder={0 0 0},
            breaklinks=true}
\urlstyle{same}  % don't use monospace font for urls
\usepackage{graphicx}
% grffile has become a legacy package: https://ctan.org/pkg/grffile
\IfFileExists{grffile.sty}{%
\usepackage{grffile}
}{}
\makeatletter
\def\maxwidth{\ifdim\Gin@nat@width>\linewidth\linewidth\else\Gin@nat@width\fi}
\def\maxheight{\ifdim\Gin@nat@height>\textheight\textheight\else\Gin@nat@height\fi}
\makeatother
% Scale images if necessary, so that they will not overflow the page
% margins by default, and it is still possible to overwrite the defaults
% using explicit options in \includegraphics[width, height, ...]{}
\setkeys{Gin}{width=\maxwidth,height=\maxheight,keepaspectratio}
\IfFileExists{parskip.sty}{%
\usepackage{parskip}
}{% else
\setlength{\parindent}{0pt}
\setlength{\parskip}{6pt plus 2pt minus 1pt}
}
\setlength{\emergencystretch}{3em}  % prevent overfull lines
\providecommand{\tightlist}{%
  \setlength{\itemsep}{0pt}\setlength{\parskip}{0pt}}
\setcounter{secnumdepth}{0}
% Redefines (sub)paragraphs to behave more like sections
\ifx\paragraph\undefined\else
\let\oldparagraph\paragraph
\renewcommand{\paragraph}[1]{\oldparagraph{#1}\mbox{}}
\fi
\ifx\subparagraph\undefined\else
\let\oldsubparagraph\subparagraph
\renewcommand{\subparagraph}[1]{\oldsubparagraph{#1}\mbox{}}
\fi

%%% Use protect on footnotes to avoid problems with footnotes in titles
\let\rmarkdownfootnote\footnote%
\def\footnote{\protect\rmarkdownfootnote}

%%% Change title format to be more compact
\usepackage{titling}

% Create subtitle command for use in maketitle
\providecommand{\subtitle}[1]{
  \posttitle{
    \begin{center}\large#1\end{center}
    }
}

\setlength{\droptitle}{-2em}

  \title{}
    \pretitle{\vspace{\droptitle}}
  \posttitle{}
    \author{}
    \preauthor{}\postauthor{}
    \date{}
    \predate{}\postdate{}
  

\begin{document}

\hypertarget{anna-oskina}{%
\section{Anna Oskina}\label{anna-oskina}}

Lecturer at HSE University, Master's student in DH.

\leavevmode\hypertarget{webaddress}{}%
annaoskina@gmail.com \textbar{} My GitHub page

\hypertarget{currently}{%
\subsection{Currently}\label{currently}}

Lecturer of contemporary and classical Japanese at HSE University

Master's 1st grade student in Digital Humanities at HSE University

\hypertarget{specialized-in}{%
\subsubsection{Specialized in}\label{specialized-in}}

Japanese classical literature, travel diaries, kuzushi-ji, xylographs,
manuscripts, book history

\hypertarget{research-interests}{%
\subsubsection{Research interests}\label{research-interests}}

DH, text mining, quantitative analyse, Meiji-Taisho fiction writings

\hypertarget{education}{%
\subsection{Education}\label{education}}

\texttt{2006-2011} \textbf{Russian State University for the Humanities,
Moscow, Russia.}

\texttt{2009-2010} \textbf{Monbukagakusho student at Chiba Univesrity,
Chiba, Japan.}

\texttt{August\ 2018} \textbf{Fifth Graduate Summer School in Japanese
Early-modern Palaeography, Emmanuel College, Cambridge.}

\texttt{February\ 2019} \textbf{EAJRS/NIJL Kuzushi-ji Workshop 2019,
Collège de France, Paris.}

\hypertarget{occupation}{%
\subsection{Occupation}\label{occupation}}

\texttt{2011-now} \textbf{Lecturer at National Research University
`Higher School of Economics', Moscow.}

https://www.hse.ru/en/staff/annaoskina

Courses:

\begin{itemize}
\tightlist
\item
  Contemporary Japanese
\item
  Classical Japanese
\item
  Japanese literature
\end{itemize}

\hypertarget{awards}{%
\subsection{Awards}\label{awards}}

\texttt{2019} Best Teacher - 2019, HSE University, Moscow.

\hypertarget{publications}{%
\subsection{Publications}\label{publications}}

\texttt{Article} Oskina A., Pasivkina S. The Precious Notes on
Papermaking» (``Kamisuki Cho:ho:ki'') as the first illustrated Guide to
Papermaking // Japanese Studies in Russia, 2018. № 3. Pp. 49-77. (in
Russian)

\texttt{Article} Oskina A. Two Diaries of the Nun Abutsu (1221?-1283):
Literature, Biography, Personality.// RGGU Bulletin. History. Philology.
Cultural Studies. Oriental Studies. 2018. № 12 (45). Pp. 251-269. (in
Russian)

\texttt{Chapter} Oskina A. Illustrated book ``Ehon Hanakazura'' //
Survival and Sustainability: Contemporary Studies in Humanities 2 Issue
314. Chiba University, 2017. Chapter. 3. Pp. 56-72. (in Russian and
Japanese)

\texttt{Chapter} Oskina A. Illustrated book ``Ehon Hanakazura'' (1765):
Education Books in the Edo Period // History and Culure of Traditional
Japan 10. Orientalia et Classica. Volume LXIX. St.Petersburg: Hyperion,
2017. Pp. 176-195. (in Russian)

\texttt{Chapter} Oskina A. Motives and Intentions of the Authors in
Diaries Based on ``The Diary of the Waning Moon'' by the Nun Abutsu and
``The Notes'' by Princess Dashkova // Literature and History:
Representation and Narrative Issue 289. Chiba University, 2015. Pp.
83-99. (in Russian)

\texttt{Article} Oskina A. ``The Diary of the Waning Moon'': Content,
Composition, Style // Vestnik Novosibirsk State University: History and
Philology. 2014. Vol. 13. № 4. Pp. 153-158. (in Russian)

\texttt{Chapter} Oskina A. Distinctive features of ``The Diary of the
Waning Moon'' by the Nun Abutsu (1221?-1283) // Oriental Societies:
Traditions and Modernity. Moscow, Baku: Center for Strategic Studies
under the President of Azerbaijan, 2014. Pp. 328-342. (in Russian)

\texttt{Chapter} Oskina A.
高等女学校の日常生活：女学生の楽しみや悩み　//　アルザス日欧知的交流事業　日本研究セミナー［大正／戦前］報告書.
Strasbourg : Centre européen d'études japonaises d'Alsace, 2014. (in
Japanese)

\texttt{Article} Oskina A. Kaibara Ekken and his ``Precepts for girls''
// Japan. Annuary. 2013. Pp. 198-215. (in Russian)

\hypertarget{conferences}{%
\subsection{Conferences}\label{conferences}}

\texttt{2019} 21th Annual Conference ``History and Culture of Japan''
(Moscow). Presentation: Travel notes of the Nun Abutsu: connection with
the literary tradition.

\texttt{2018} 20th Annual Conference ``History and Culture of Japan''
(Moscow). Presentation: The History of the Reizei Family Library.

\texttt{2017} 19th Annual Conference ``History and Culture of Japan''
(Moscow). Presentation: Education with a book in the Edo Period based on
``Ehon Hanakazura'' (1765).

\texttt{2014} 16th Annual Conference ``History and Culture of Japan''
(Moscow). Presentation: Prose and poetry in the nun Abutsu's
(1221?-1283) ``The Diary of the Waning Moon''.

\texttt{2013} 15th Annual Conference ``History and Culture of Japan''
(Moscow). Presentation: The nun Abutsu and ``The Tale of Genji''.

Written Monuments of the East. The Problems of Interpretation and
Translation (Moscow). Presentation: The problems in translation of
``Izayoi nikki'' by the nun Abutsu: quotes and allusions.

European Association of Japanese Resource Specialists Conference 2013
(Paris). Presentation: Illustrated book ``Ehon hanakazura'' (1765):
pictures, poems, proverbs.

2nd Congress of Young Orientalists of Russia and the CIS (Baku).
Presentation: Distinctive features of ``The Diary of the Waning Moon''
by the nun Abutsu (1221?-1283).

\texttt{2012} 14th Annual Conference ``History and Culture of Japan''
(Moscow). Presentation: From capital to capital: the travelling in Heian
and Edo periods.

Nikolai Nevsky Memorial Conference - International Symposium in Honor of
the 120th Birthday of N.A. Nevsky (St.~Petersburg). Presentation: The
diaries of the nun Abutsu as a source for her biography.

\texttt{2011} 13th Annual Conference ``History and Culture of Japan''
(Moscow). Presentation: ``The Great Learning for women'': from Kaibara
Ekken to Fukuzawa Yukichi.

Japanese Study Seminar - Taisho Prewar (Showa) (Alsace). Presentation:
Everyday Life in Women High Schools in the Taisho period.

\texttt{2010} Conference of Young Japanologists (St.~Petersburg).
Presentation: Kaibara Ekken and Fukuzawa Yukichi: old and new views on
women education.

\texttt{2009} 11th Annual Conference ``History and Culture of Japan''
(Moscow). Presentation: Kaibara Ekken and women learning in the Edo
period.


\end{document}
